%!TEX root = forallxyyc.tex
\part{Noções fundamentais de lógica}
\label{ch.intro}
\addtocontents{toc}{\protect\mbox{}\protect\hrulefill\par}


\chapter{Argumentos}
\label{s:Arguments}

A lógica tem \emph{muitos} usos, como mencionado no~\hyperref[preface]{prefácio}. Aqui, vamos nos concentrar em seu uso na avaliação de argumentos: separar os bons dos ruins.

Na linguagem cotidiana, às vezes usamos a palavra \textquotedblleft argumento\textquotedblright{} para falar de discussões acaloradas aos berros. Se você e um amigo têm um argumento nesse sentido, as coisas não vão bem entre vocês. A lógica não se ocupa desse ranger de dentes e puxar de cabelos. Isso não são argumentos no sentido em que usamos o termo; são apenas desentendimentos.

Um argumento, tal como o entenderemos, é mais parecido com isto:
	\begin{earg}\label{argButlerGardner}
		\item Ou foi o mordomo ou foi o jardineiro.
		\item O mordomo não fez isso.
		\item[\texttherefore] O jardineiro fez isso.
	\end{earg}
Temos aqui uma série de sentenças. Os três pontos na terceira linha do argumento são lidos como \textquotedblleft portanto\textquotedblright{}. Eles indicam que a sentença final expressa a \emph{conclusão} do argumento. As duas sentenças anteriores são as \emph{premissas} do argumento. Se você acredita nas premissas e considera que a conclusão decorre das premissas — isto é, que o argumento, como diremos, é válido —, então isso talvez lhe dê uma razão para acreditar na conclusão.

É esse tipo de coisa que interessa aos lógicos. Diremos que um argumento é qualquer conjunto de premissas acompanhado de uma conclusão.

Esta parte discute algumas noções lógicas básicas aplicáveis a argumentos em uma língua natural, como o português. É importante começar com uma compreensão clara do que são os argumentos e do que significa um argumento ser válido. Mais adiante, representaremos argumentos em português em uma linguagem formal. %Queremos que a validade formal, tal como definida na linguagem formal, tenha ao menos algumas das características importantes da validade em língua natural.

No exemplo que acabamos de dar, usamos sentenças individuais para expressar tanto as duas premissas do argumento quanto uma terceira sentença para expressar sua conclusão. Muitos argumentos são expressos dessa maneira, mas uma única sentença pode conter um argumento completo. Considere:
	\begin{quote}
		 O mordomo tem um álibi; portanto, não pode ter feito isso.
	\end{quote}
Esse argumento tem uma premissa seguida de uma conclusão.

Muitos argumentos começam com premissas e terminam com uma conclusão, mas não todos. O argumento com o qual esta seção começou poderia igualmente ter sido apresentado com a conclusão no início, assim:
	\begin{quote}
		O jardineiro fez isso. Afinal, foi ou o mordomo ou o
		jardineiro. E o mordomo não fez isso.
	\end{quote}
Da mesma forma, ele poderia ter sido apresentado com a conclusão no meio:
	\begin{quote}
		O mordomo não fez isso. Consequentemente, foi o jardineiro,
		dado que foi ou o jardineiro ou o mordomo.
	\end{quote}
Ao abordar um argumento, queremos saber se a conclusão decorre ou não das premissas. Portanto, a primeira coisa a fazer é separar a conclusão das premissas. Como orientação, estas palavras são frequentemente usadas para indicar a conclusão de um argumento:
	\begin{center}
		portanto, logo, daí, assim, consequentemente, por conseguinte
	\end{center}
Por esse motivo, às vezes são chamadas de \define{palavras indicadoras de conclus\~ao}.

Em contraste, estas expressões são \define{palavras indicadoras de premissa},
pois frequentemente indicam que estamos diante de uma premissa, e não de uma
conclusão:
	\begin{center}
		como, porque, dado que
	\end{center}
Mas, ao analisar um argumento, não há substituto para um bom faro.

\newglossaryentry{premise indicator word}
{
name=indicador de premissa,
description={Uma palavra ou expressão, como \textquotedblleft porque\textquotedblright{}, usada para indicar que o que vem a seguir é a premissa de um argumento}
}

\newglossaryentry{conclusion indicator word}
{
name=indicador de conclusão,
description={Uma palavra ou expressão, como \textquotedblleft portanto\textquotedblright{}, usada para indicar que o que vem a seguir é a conclusão de um argumento}
}

\newglossaryentry{argument}
{
name=argumento,
description={Uma série encadeada de sentenças, dividida em \gls{premise}s e \gls{conclusion}}
}

\newglossaryentry{premise}
{
name=premissa,
description={Uma sentença em um \gls{argument} que não é a \gls{conclusion}}
}

\newglossaryentry{conclusion}
{
name=conclusão,
description={A última sentença em um \gls{argument}}
}


\section{Sentenças}
\label{intro.sentences}

Para sermos perfeitamente gerais, podemos definir um \define{argumento} como uma série de sentenças. As sentenças no início da série são premissas. A sentença final da série é a conclusão. Se as premissas forem verdadeiras e o argumento for bom, então você tem uma razão para aceitar a conclusão.

Na lógica, estamos interessados apenas em sentenças que possam figurar como premissa ou conclusão de um argumento, isto é, sentenças que possam ser verdadeiras ou falsas. Assim, restringiremos nossa atenção a sentenças desse tipo e definiremos uma \define{senten\c{c}a} como uma sentença que pode ser verdadeira ou falsa.

Você não deve confundir a ideia de uma sentença que pode ser verdadeira ou falsa com a diferença entre fato e opinião. Muitas vezes, as sentenças da lógica expressarão coisas que seriam consideradas fatos — como \textquotedblleft Rudolf Carnap nasceu em Ronsdorf\textquotedblright{} ou \textquotedblleft Simone de Beauvoir gostava de fazer caminhadas\textquotedblright{}. Elas também podem expressar coisas que você talvez considere questões de opinião — como \textquotedblleft O ruibarbo é saboroso\textquotedblright{}. Em outras palavras, uma sentença não é desqualificada para fazer parte de um argumento porque não sabemos se é verdadeira ou falsa, ou porque sua verdade ou falsidade é uma questão de opinião. Se for o tipo de sentença que pode ser verdadeira ou falsa, ela pode desempenhar o papel de premissa ou conclusão.

Além disso, há coisas que seriam consideradas \textquotedblleft sentenças\textquotedblright{} em um curso de linguística ou gramática, mas que não consideraremos sentenças na lógica.

\paragraph{Perguntas} Em uma aula de gramática, \textquotedblleft Você ainda está com sono?\textquotedblright{} seria uma sentença interrogativa. Embora você possa estar com sono ou alerta, a pergunta em si não é verdadeira nem falsa. Por esse motivo, perguntas não contarão como sentenças na lógica. Suponha que você responda à pergunta: \textquotedblleft Não estou com sono.\textquotedblright{} Isso é verdadeiro ou falso e, portanto, é uma sentença no sentido lógico. Em geral, \emph{perguntas} não contarão como sentenças, mas \emph{respostas} contarão.

\textquotedblleft Sobre o que é este curso?\textquotedblright{} não é uma sentença (no nosso sentido). \textquotedblleft Ninguém sabe sobre o que é este curso.\textquotedblright{} é uma sentença.

\paragraph{Imperativos} Comandos são frequentemente formulados como imperativos, como \textquotedblleft Acorde!\textquotedblright{}, \textquotedblleft Sente-se direito\textquotedblright{} e assim por diante. Em uma aula de gramática, eles seriam considerados sentenças imperativas. Embora possa ser bom para você sentar-se direito ou talvez não, o comando não é verdadeiro nem falso. Observe, contudo, que os comandos nem sempre são formulados como imperativos. \textquotedblleft Você respeitará minha autoridade\textquotedblright{} \emph{é} verdadeiro ou falso — ou você a respeitará ou não — e, portanto, conta como uma sentença no sentido lógico.

\paragraph{Exclamações} \textquotedblleft Ai!\textquotedblright{} às vezes é chamada de sentença exclamativa, mas não é verdadeira nem falsa. Trataremos \textquotedblleft Ai, machuquei o dedo do pé!\textquotedblright{} como significando o mesmo que \textquotedblleft Machuquei o dedo do pé.\textquotedblright{} O \textquotedblleft ai\textquotedblright{} não acrescenta nada que possa ser verdadeiro ou falso.


\practiceproblems
No final de alguns capítulos, há exercícios que revisam e exploram o material abordado no capítulo. Não há substituto para resolver efetivamente alguns problemas, pois aprender lógica consiste mais em desenvolver uma maneira de pensar do que em memorizar fatos.

\medskip

Eis o primeiro exercício. Destaque a expressão que exprime a conclusão de cada um destes argumentos:
\begin{compactlist}
	\item Está ensolarado. Portanto, devo levar meus óculos de sol.
	\item Deve ter feito sol. Afinal, eu usei meus óculos de sol.
	\item Ninguém além de você pôs as mãos no pote de biscoitos. E a cena do crime está coberta de migalhas de biscoito. Você é o culpado!
	\item A senhorita Scarlett e o professor Plum estavam na sala de estudo no
	momento do assassinato. O reverendo Green estava com o castiçal no
	salão de baile, e sabemos que ele não tem sangue nas mãos. Logo, o
	coronel Mustard fez isso na cozinha com o cano de chumbo.
	Afinal, a arma não havia sido disparada.
\end{compactlist}


\chapter{O escopo da lógica}
\label{s:Valid}

\section{Consequência e validade}

Em \cref{s:Arguments}, falamos de argumentos, isto é, de conjuntos de sentenças (as premissas), seguidas por uma única sentença (a conclusão). Dissemos que algumas palavras, como \textquotedblleft portanto\textquotedblright{}, indicam qual sentença deve ser a conclusão. \textquotedblleft Portanto\textquotedblright{}, é claro, sugere que há uma conexão entre as premissas e a conclusão, a saber, que a conclusão \emph{decorre das} premissas ou \emph{é uma consequência das} premissas.

Essa noção de consequência é uma das principais preocupações da lógica. Poderíamos até dizer que a lógica é a ciência do que decorre de quê. A lógica desenvolve teorias e ferramentas que nos dizem quando uma sentença decorre de outras.

E quanto ao argumento principal discutido em \cref{s:Arguments}? 
\begin{earg}
	\item[] Ou foi o mordomo ou foi o jardineiro.
	\item[] O mordomo não fez isso.
	\item[\texttherefore] O jardineiro fez isso.
\end{earg}
Não temos nenhum contexto que nos diga a que se referem as sentenças desse argumento. Talvez você suspeite que \textquotedblleft fez isso\textquotedblright{} signifique aqui \textquotedblleft foi o autor\textquotedblright{} de algum crime não especificado. Você pode imaginar que o argumento ocorre em um romance policial ou em um programa de televisão, talvez dito por um detetive que examina as evidências. Mas, mesmo sem nenhuma dessas informações, você provavelmente concorda que o argumento é bom no sentido de que, seja lá ao que as premissas se refiram, se ambas forem verdadeiras, a conclusão também não poderá deixar de ser verdadeira. Se a primeira premissa for verdadeira, isto é, se for verdade que \textquotedblleft o mordomo fez isso ou o jardineiro fez isso\textquotedblright{}, então pelo menos um deles \textquotedblleft fez isso\textquotedblright{}, seja lá o que isso signifique. E, se a segunda premissa for verdadeira, então o mordomo não \textquotedblleft fez isso\textquotedblright{}. Resta apenas uma opção: \textquotedblleft o jardineiro fez isso\textquotedblright{} deve ser verdadeiro. Aqui, a conclusão decorre das premissas. Chamamos de \define{válidos} os argumentos que têm essa propriedade.

Em contraste, considere o seguinte argumento:
\begin{earg}\label{argMaidDriver}
	\item[] Se o motorista fez isso, a criada não fez isso.
	\item[] A criada não fez isso.
	\item[\texttherefore] O motorista fez isso.
\end{earg}
Ainda não temos ideia do que está sendo discutido aqui. Mas, novamente, você provavelmente concorda que esse argumento é diferente do anterior em um aspecto importante. Se as premissas forem verdadeiras, não é garantido que a conclusão também seja verdadeira. As premissas desse argumento não excluem, por si sós, que alguém que não seja a criada nem o motorista tenha feito isso. Portanto, há um caso em que ambas as premissas são verdadeiras e, ainda assim, o motorista não fez isso, isto é, a conclusão não é verdadeira. Nesse segundo argumento, a conclusão não decorre das premissas. Se, como nesse argumento, a conclusão não decorre das premissas, dizemos que ele é \define{inválido}.

\section{Casos e tipos de validade}

Como determinamos que o segundo argumento é inválido? Apontamos para um caso em que as premissas são verdadeiras e a conclusão não é. Esse era o cenário em que nem o motorista nem a criada fizeram isso, mas uma terceira pessoa fez. Chamemos esse caso de \define{contraexemplo} do argumento. Se há um contraexemplo para um argumento, a conclusão não pode ser uma consequência das premissas. Para que a conclusão seja uma consequência das premissas, a verdade das premissas deve garantir a verdade da conclusão. Se há um contraexemplo, a verdade das premissas não garante a verdade da conclusão.

Como lógicos, queremos ser capazes de determinar quando a conclusão de um argumento decorre das premissas. E a conclusão é uma consequência das premissas se não há contraexemplo — nenhum caso em que todas as premissas sejam verdadeiras, mas a conclusão não seja. Isso motiva uma definição:

	\factoidbox{
		Uma sentença $A$ é uma \define{consequência} das sentenças $B_1$, \dots, $B_n$ se, e somente se, não há caso em que $B_1$, \dots, $B_n$ sejam todas verdadeiras e $A$ não seja verdadeira. (Também dizemos que $A$ \define{decorre de} $B_1$, \dots, $B_n$ ou que $B_1$, \dots, $B_n$ \define{implicam}~$A$.)
	}

Essa \textquotedblleft definição\textquotedblright{} é incompleta: ela não nos diz o que é um \textquotedblleft caso\textquotedblright{} nem o que significa ser \textquotedblleft verdadeiro em um caso\textquotedblright{}. Até agora, vimos apenas um exemplo: um cenário hipotético envolvendo três pessoas. Das três pessoas no cenário — um motorista, uma criada e uma terceira pessoa qualquer —, o motorista e a criada não fizeram isso, mas a terceira pessoa fez. Nesse cenário, tal como descrito, o motorista não fez isso e, portanto, é um caso em que a sentença \textquotedblleft o motorista fez isso\textquotedblright{} não é verdadeira. As premissas do nosso segundo argumento são verdadeiras, mas a conclusão não é verdadeira: o cenário é um contraexemplo.

Dissemos que os argumentos em que a conclusão é uma consequência das premissas são chamados de válidos, e aqueles em que a conclusão não é uma consequência das premissas são inválidos. Como agora temos pelo menos uma primeira tentativa de definição de \textquotedblleft consequência\textquotedblright{}, vamos registrar isto:

	\factoidbox{
		\begin{factoiditems}
			\item Um argumento é \define{válido} se, e somente se, a conclusão é uma consequência das premissas.
			\item Um argumento é \define{inválido} se, e somente se, não é válido, isto é, tem um contraexemplo.
		\end{factoiditems}}

\newglossaryentry{valid}
{
name=válido,
description={Uma propriedade de argumentos cuja conclusão é uma consequência das premissas}
}

\newglossaryentry{invalid}
{
name=inválido,
description={Uma propriedade de argumentos que ocorre quando a conclusão não é uma consequência das premissas; o oposto de \gls{valid}}
}

O trabalho dos lógicos é tornar mais precisa a noção de \textquotedblleft caso\textquotedblright{} e investigar quais argumentos são válidos quando \textquotedblleft caso\textquotedblright{} é precisado de uma maneira ou de outra. Se tomarmos \textquotedblleft caso\textquotedblright{} como significando \textquotedblleft cenário hipotético\textquotedblright{}, como no contraexemplo ao segundo argumento, fica claro que o primeiro argumento conta como válido. Se imaginarmos um cenário em que ou o mordomo ou o jardineiro fez isso e em que o mordomo também não fez isso, estaremos automaticamente imaginando um cenário em que o jardineiro fez isso. Assim, qualquer cenário hipotético em que as premissas do nosso primeiro argumento sejam verdadeiras automaticamente torna verdadeira a conclusão do nosso primeiro argumento. Isso torna o primeiro argumento válido.

Tornar \textquotedblleft caso\textquotedblright{} mais específico, interpretando-o como \textquotedblleft cenário hipotético\textquotedblright{}, é um avanço. Mas não é o fim da história. O primeiro problema é que não sabemos o que deve ser considerado um cenário hipotético. Eles são limitados pelas leis da física? Pelo que é concebível, em um sentido muito geral? As respostas que dermos a essas perguntas determinam quais argumentos consideramos válidos.

Suponha que a resposta à primeira pergunta seja \textquotedblleft sim\textquotedblright{}. Considere o seguinte argumento:
	\begin{earg}
		\item[] A espaçonave \textit{Rocinante} levou seis horas para chegar a Júpiter a partir da estação espacial Tycho.
		\item[\texttherefore] A distância entre a estação espacial Tycho e Júpiter é inferior a 14~bilhões de quilômetros.
	\end{earg}
Um contraexemplo a esse argumento seria um cenário em que a \textit{Rocinante} percorresse mais de 14 bilhões de quilômetros em 6 horas, ultrapassando a velocidade da luz. Como tal cenário é incompatível com as leis da física, não há um cenário assim se os cenários hipotéticos tiverem de obedecer às leis da física. Se os cenários hipotéticos não forem limitados pelas leis da física, contudo, há um contraexemplo: um cenário em que a \textit{Rocinante} viaja mais rápido que a luz.

Suponha que a resposta à segunda pergunta seja \textquotedblleft sim\textquotedblright{} e considere outro argumento:
	\begin{earg}
		\item[] Priya é uma oftalmologista.
		\item[\texttherefore] Priya é uma médica dos olhos.
	\end{earg}
Se permitirmos apenas cenários concebíveis, este também é um argumento válido. Ao imaginar Priya sendo uma oftalmologista, você imagina, com isso, Priya sendo uma médica dos olhos. É isso que \textquotedblleft oftalmologista\textquotedblright{} e \textquotedblleft médica dos olhos\textquotedblright{} significam. Um cenário em que Priya é uma oftalmologista, mas não uma médica dos olhos, é excluído pela conexão conceitual entre essas palavras.

Dependendo dos tipos de casos que consideramos contraexemplos potenciais, chegamos a diferentes noções de consequência e validade. A validade é caracterizada pela ausência de contraexemplos. Mas o que conta como contraexemplo pode ser diferente para diferentes noções. Por exemplo, podemos excluir contraexemplos que violam as leis da natureza ou contraexemplos que violam conexões conceituais entre palavras como \textquotedblleft médica dos olhos\textquotedblright{} e \textquotedblleft oftalmologista\textquotedblright{}. Assim, podemos chamar um argumento de \define{nomologicamente válido} se não há contraexemplos que respeitem as leis da natureza, e de \define{conceitualmente válido} se não há contraexemplos que respeitem as conexões conceituais entre palavras. Para ambas essas noções de validade, aspectos do mundo (por exemplo, quais são as leis da natureza) e aspectos do significado das sentenças no argumento (por exemplo, que \textquotedblleft oftalmologista\textquotedblright{} significa simplesmente um tipo de médica dos olhos) contribuem para determinar se um argumento conta como válido.

\section{Validade formal}\label{s:FormalValidity}

Uma característica que distingue a \emph{consequência lógica}, contudo, é que ela não deve depender do conteúdo das premissas e da conclusão, mas apenas de sua forma lógica. Em outras palavras, como lógicos, queremos desenvolver uma teoria que capte outra maneira pela qual os argumentos podem ser válidos. Por exemplo, tanto
\begin{earg}
	\item[] Priya é uma oftalmologista ou uma dentista.
	\item[] Priya não é uma dentista.
	\item[\texttherefore] Priya é uma médica dos olhos.
\end{earg}
quanto
\begin{earg}
	\item[] Priya é uma oftalmologista ou uma dentista.
	\item[] Priya não é uma dentista.
	\item[\texttherefore] Priya é uma oftalmologista.
\end{earg}
são argumentos válidos. Mas, embora a validade do primeiro dependa do conteúdo (isto é, do significado de \textquotedblleft oftalmologista\textquotedblright{} e \textquotedblleft médica dos olhos\textquotedblright{}), a do segundo não depende. O segundo argumento é \define{formalmente válido}. Podemos descrever a \textquotedblleft forma\textquotedblright{} desse argumento como um padrão, algo como:
\begin{earg}
	\item[] $A$ é uma $X$ ou uma $Y$.
	\item[] $A$ não é uma $Y$.
	\item[\texttherefore] $A$ é uma $X$.
\end{earg}
Aqui, $A$, $X$ e $Y$ são marcadores de posição para expressões apropriadas que, quando substituídas por $A$, $X$ e $Y$, transformam o padrão em um argumento constituído por sentenças. Por exemplo,
\begin{earg}
	\item[] Mei é uma matemática ou uma botânica.
	\item[] Mei não é uma botânica.
	\item[\texttherefore] Mei é uma matemática.
\end{earg}
é um argumento da mesma forma, mas o primeiro argumento acima não é: teríamos de substituir $Y$ por expressões diferentes (uma vez por \textquotedblleft oftalmologista\textquotedblright{} e outra vez por \textquotedblleft médica dos olhos\textquotedblright{}) para obtê-lo a partir do padrão.

Além disso, o primeiro argumento não é formalmente válido. \emph{Sua} forma é:
\begin{earg}
	\item[] $A$ é uma $X$ ou uma $Y$.
	\item[] $A$ não é uma $Y$.
	\item[\texttherefore] $A$ é uma $Z$.
\end{earg}
Nesse padrão, podemos substituir $X$ por \textquotedblleft oftalmologista\textquotedblright{} e $Z$ por \textquotedblleft médica dos olhos\textquotedblright{} para obter o argumento original. Mas eis outro argumento da mesma forma:
\begin{earg}
	\item[] Mei é uma matemática ou uma botânica.
	\item[] Mei não é uma botânica.
	\item[\texttherefore] Mei é uma acrobata.
\end{earg}
Esse argumento claramente não é válido, pois podemos imaginar uma matemática chamada Mei que não seja uma acrobata.

Nossa estratégia como lógicos será propor uma noção de \textquotedblleft caso\textquotedblright{} segundo a qual um argumento só será válido se for formalmente válido. É claro que tal noção de \textquotedblleft caso\textquotedblright{} terá de violar não apenas algumas leis da natureza, mas também algumas leis do português. Como o primeiro argumento é inválido nesse sentido, devemos permitir como contraexemplo um caso em que Priya seja uma oftalmologista, mas não uma médica dos olhos. Essa situação não é concebível: ela é excluída pelos significados de \textquotedblleft oftalmologista\textquotedblright{} e \textquotedblleft médica dos olhos\textquotedblright{}.

Quando considerarmos casos de vários tipos para avaliar a validade de um argumento, faremos uma suposição importante. Supomos que todo caso torna toda sentença sob consideração verdadeira ou não verdadeira. Isso significa, antes de tudo, que qualquer cenário imaginado que deixe indeterminado se uma sentença do nosso argumento é verdadeira não será considerado um contraexemplo potencial. Por exemplo, um cenário em que Priya é uma dentista, mas não uma oftalmologista, contará como um caso a ser considerado nos primeiros argumentos desta seção, mas não como um caso a ser considerado nos dois últimos: ele não nos diz se Mei é uma matemática, uma botânica ou uma acrobata. Se um caso não torna uma sentença verdadeira, dizemos que ele a torna \define{falsa}. Assim, suporemos que os casos tornam as sentenças verdadeiras ou falsas, mas nunca ambas.\footnote{A suposição de que as sentenças são verdadeiras ou falsas, e não ambas, é chamada de \textquotedblleft bivalência\textquotedblright{}. Mesmo que essa suposição pareça fazer sentido comum, ela é controversa entre filósofos da lógica. Primeiro, há lógicos que querem considerar casos em que as sentenças não são nem verdadeiras nem falsas, mas têm algum tipo de nível intermediário de verdade. De maneira ainda mais controversa, alguns filósofos pensam que deveríamos permitir a possibilidade de as sentenças serem verdadeiras e falsas ao mesmo tempo. Há sistemas de lógica nos quais as sentenças podem não ser nem verdadeiras nem falsas, ou podem ser ambas, mas não discutiremos esses sistemas neste livro.}

\section{Argumentos corretos}

Antes de prosseguirmos e executarmos essa estratégia, são necessárias algumas explicações. Definimos argumentos como conjuntos de sentenças em que uma delas (a conclusão) deve decorrer das outras (as premissas). Argumentos nesse sentido são usados o tempo todo no discurso cotidiano e científico. Quando são usados, os argumentos são apresentados para sustentar ou até provar suas conclusões. Ora, se um argumento é válido, ele sustentará sua conclusão, mas \emph{somente se} todas as suas premissas forem verdadeiras. A validade exclui a possibilidade de que as premissas sejam verdadeiras e a conclusão não seja verdadeira ao mesmo tempo. Ela não exclui, por si só, a possibilidade de que a conclusão não seja verdadeira, ponto final. Em outras palavras, é perfeitamente possível que um argumento válido tenha uma conclusão que não seja verdadeira!

Considere este exemplo:
	\begin{earg}
		\item[] Laranjas são frutas ou instrumentos musicais.
		\item[] Laranjas não são frutas.
		\item[\texttherefore] Laranjas são instrumentos musicais.
	\end{earg}
A conclusão desse argumento é ridícula. No entanto, ela decorre das premissas. \emph{Se} ambas as premissas forem verdadeiras, \emph{então} a conclusão terá necessariamente de ser verdadeira. Portanto, o argumento é válido.

Por outro lado, ter premissas verdadeiras e uma conclusão verdadeira não é suficiente para tornar um argumento válido. Considere este exemplo:
	\begin{earg}
		\item[] Londres fica na Inglaterra.
		\item[] Pequim fica na China.
		\item[\texttherefore] Paris fica na França.
	\end{earg}
As premissas e a conclusão desse argumento são, de fato, todas verdadeiras, mas o argumento é inválido. Se Paris declarasse independência do restante da França, a conclusão deixaria de ser verdadeira, embora as duas premissas continuassem verdadeiras. Assim, há um caso em que as premissas desse argumento são verdadeiras sem que a conclusão seja verdadeira. Portanto, o argumento é inválido.

O importante é lembrar que a validade não diz respeito à verdade ou falsidade \emph{real} das sentenças no argumento. Diz respeito a saber se é \emph{possível} que todas as premissas sejam verdadeiras e a conclusão não seja verdadeira ao mesmo tempo (em algum caso hipotético). O que \emph{de fato} ocorre não tem nenhum papel especial; e quais sejam os fatos não determina se um argumento é válido ou não.\footnote{Bem, há um caso em que determina: se as premissas forem de fato verdadeiras e a conclusão for de fato não verdadeira, então vivemos em um contraexemplo; logo, o argumento é inválido.} Nada sobre o modo como as coisas são pode, por si só, determinar se um argumento é válido. Costuma-se dizer que a lógica não se importa com sentimentos. Na verdade, ela também não se importa com fatos.

Quando usamos um argumento para provar que sua conclusão \emph{é verdadeira}, precisamos, então, de duas coisas. Primeiro, precisamos que o argumento seja válido; isto é, precisamos que a conclusão decorra das premissas. Mas também precisamos que as premissas sejam verdadeiras. Diremos que um argumento é \define{correto} se, e somente se, é válido e todas as suas premissas são verdadeiras.

\newglossaryentry{sound}
{
name=correto,
description={Uma propriedade de argumentos que ocorre quando o argumento é válido e tem todas as premissas verdadeiras}
}

O outro lado da questão é que, quando você quer refutar um argumento, tem duas opções: pode mostrar que uma ou mais premissas não são verdadeiras ou pode mostrar que o argumento não é válido. A lógica, contudo, só ajudará você com a segunda opção!

\section{Argumentos indutivos}

Muitos bons argumentos são inválidos. Considere este:
	\begin{earg}
		\item[] Em todos os invernos até agora, nevou em Calgary.
	\item[\texttherefore] Nevará em Calgary no próximo inverno.
\end{earg}
Esse argumento generaliza observações sobre muitos casos (passados) para uma conclusão sobre todos os casos (futuros). Esses argumentos são chamados de argumentos \define{indutivos}. Ainda assim, o argumento é inválido. Mesmo que tenha nevado em Calgary em todos os invernos até agora, continua \emph{possível} que Calgary permaneça seca durante todo o próximo inverno. Na verdade, mesmo que daqui em diante neve todo inverno em Calgary, ainda poderíamos \emph{imaginar} um caso em que este ano seja o primeiro ano em que não neve durante todo o inverno. E esse cenário hipotético é um caso em que as premissas do argumento são verdadeiras, mas a conclusão não é, tornando o argumento inválido.

O ponto de tudo isso é que argumentos indutivos — mesmo bons argumentos indutivos — não são válidos (dedutivamente). Eles não são \emph{à prova de falhas}. Por improvável que seja, é \emph{possível} que sua conclusão seja falsa, mesmo quando todas as suas premissas são verdadeiras. Neste livro, deixaremos de lado (inteiramente) a questão do que torna bom um argumento indutivo. Nosso interesse é simplesmente separar os argumentos (dedutivamente) válidos dos inválidos.

Assim, estamos interessados em saber se uma conclusão \emph{decorre de} algumas premissas ou não. Não diga, contudo, que as premissas \emph{inferem} a conclusão. A consequência lógica é uma relação entre premissas e conclusões; a inferência é algo que fazemos. Portanto, se quiser mencionar a inferência quando a conclusão decorre das premissas, você pode dizer que \emph{se pode inferir} a conclusão das premissas.


\practiceproblems
\problempart
Quais dos seguintes argumentos são válidos? Quais são inválidos?
\begin{compactlist}
\item
\begin{earg}
\item Sócrates é um homem.
\item Todos os homens são cenouras.
\item[\texttherefore] Sócrates é uma cenoura.
\end{earg}
\item
\begin{earg}
\item Abe Lincoln nasceu em Illinois ou alguma vez foi presidente.
\item Abe Lincoln nunca foi presidente.
\item[\texttherefore] Abe Lincoln nasceu em Illinois.
\end{earg}
\item
\begin{earg}
\item Se eu puxar o gatilho, Abe Lincoln morrerá.
\item Eu não puxo o gatilho.
\item[\texttherefore] Abe Lincoln não morrerá.
\end{earg}
\item
\begin{earg}
\item Abe Lincoln era da França ou de Luxemburgo.
\item Abe Lincoln não era de Luxemburgo.
\item[\texttherefore] Abe Lincoln era da França.
\end{earg}
\item
\begin{earg}
\item Se o mundo acabar hoje, não precisarei me levantar amanhã de manhã.
\item Precisarei me levantar amanhã de manhã.
\item[\texttherefore] O mundo não acabará hoje.
\end{earg}
\item
\begin{earg}
\item Joe tem agora 19 anos.
\item Joe tem agora 87 anos.
\item[\texttherefore] Bob tem agora 20 anos.
\end{earg}
\end{compactlist}

\problempart
\label{pr.EnglishCombinations}
Poderia existir:
	\begin{compactlist}
		\item Um argumento válido que tenha uma premissa falsa e uma premissa verdadeira?
		\item Um argumento válido que tenha apenas premissas falsas?
		\item Um argumento válido com apenas premissas falsas e uma conclusão falsa?
		\item Um argumento inválido que possa se tornar válido com a adição de uma nova premissa?
		\item Um argumento válido que possa se tornar inválido com a adição de uma nova premissa?
	\end{compactlist}
Em cada caso: se sim, dê um exemplo; se não, explique por quê.

\chapter{Other logical notions}\label{s:BasicNotions}

In \cref{s:Valid}, we introduced the idea of consequence, i.e., of
validity of arguments. This is one of the most important ideas in
logic. In this section, we will introduce some similarly important
ideas. They all rely, as did validity, on the idea that sentences are
true (or not) in cases. For the rest of this section, we'll take cases
in the sense of conceivable scenario, i.e., in the sense in which we
used them to define conceptual validity. The points we made about
different kinds of validity can be made about our new notions along
similar lines: if we use a different idea of what counts as a ``case''
we will get different notions.  And as logicians we will, eventually,
consider a more permissive definition of case than we do here.

%\section{Truth values}
%As we said in \cref{s:Arguments}, arguments consist of premises and a conclusion. Note that many kinds of English sentence cannot be used to express premises or conclusions of arguments. For example:
%	\begin{itemize}
%		\item \textbf{Questions}, e.g.\ `are you feeling sleepy?'
%		\item \textbf{Imperatives}, e.g.\ `Wake up!'
%		\item \textbf{Exclamations}, e.g.\ `Ouch!'
%	\end{itemize}
%The common feature of these three kinds of sentence is that they are not \emph{assertoric}: they cannot be true or false. It does not even make sense to ask whether a \emph{question} is true (it only makes sense to ask whether the \emph{answer} to a question is true).

%The general point is that, the premises and conclusion of an argument must be capable of having a \define{truth value}. The two truth values that concern us are just True and False.

\section{Joint possibility}

Consider these two sentences:
	\begin{compactlist}
		\item[B1.] Jane's only brother is shorter than her.
		\item[B2.] Jane's only brother is taller than her.
	\end{compactlist}
Logic alone cannot tell us which, if either, of these sentences is true. Yet we can say that \emph{if} the first sentence (B1) is true, \emph{then} the second sentence (B2) must be false. Similarly, if B2 is true, then B1 must be false. There is no possible scenario where both sentences are true together. These sentences are incompatible with each other, they cannot all be true at the same time. This motivates the following definition:
	\factoidbox{
		Sentences are \define{jointly possible} if and only if there is a case where they are all true together.
	}
B1 and B2 are \emph{jointly impossible}, while, say, the following two sentences are jointly possible:
	\begin{compactlist}
		\item[B1.] Jane's only brother is shorter than her.
		\item[B3.] Jane's only brother is younger than her.
	\end{compactlist}

\newglossaryentry{possibility}
{
name=joint possibility,
text={jointly possible},
description={A property possessed by some sentences when they are all true in a single case}
}

We can ask about the joint possibility of any number of sentences. For example, consider the following four sentences:
	\begin{compactlist}	
		\item[G1.] \label{MartianGiraffes} There are at least four giraffes at the wild animal park.
		\item[G2.] There are exactly seven gorillas at the wild animal park.
		\item[G3.] There are not more than two Martians at the wild animal park.
		\item[G4.] Every giraffe at the wild animal park is a Martian.
	\end{compactlist}
G1 and G4 together entail that there are at least four Martian
giraffes at the park. This conflicts with G3, which implies that there
are no more than two Martian giraffes there. So the sentences G1--G4
are jointly impossible. They cannot all be true together. (Note that
the sentences G1, G3 and G4 are jointly impossible. But if sentences
are already jointly impossible, adding an extra sentence to the mix
cannot make them jointly possible!)

There is one thing worth pointing out: You might think that an
argument only ``makes sense'' if its premises are jointly possible.
But neither our definition of what an argument is, nor of when it is
valid, requires this. In fact, according to our definition, any
argument with jointly impossible premises is automatically valid!
(Exercise: convince yourself that this is true.)

\ifHTMLtarget
\section{Necessary truths, necessary falsehoods, and contingency}
\else
\section[Necessary truths and falsehoods]{Necessary truths, necessary falsehoods, and contingency}
\fi

In assessing arguments for validity, we care about what would be true \emph{if} the premises were true, but some sentences just \emph{must} be true. Consider these sentences:
	\begin{enumerate}
		\item\label{Acontingent} It is raining.
		\item\label{Atautology} Either it is raining here, or it is not.
		\item\label{Acontradiction} It is both raining here and not raining here.
	\end{enumerate}
In order to know if \cref*{Acontingent} is true, you would need to look outside or check the weather channel. It might be true; it might be false. A sentence which is capable of being true and capable of being false (in different circumstances, of course) is called \define{contingent}.

\newglossaryentry{contingent sentence}
{
name=contingent sentence,
description={A sentence that is neither a \gls{necessary truth} nor a \gls{necessary falsehood}; a sentence that in some case is true and in some other case, false}
}

\Cref*{Atautology} is different. You do not need to look outside to know that it is true. Regardless of what the weather is like, it is either raining or it is not. That is a \define{necessary truth}.

\newglossaryentry{necessary truth}
{
name={necessary truth},
description={A sentence that is true in every case}
}

Equally, you do not need to check the weather to determine whether or not \cref*{Acontradiction} is true. It must be false, simply as a matter of logic. It might be raining here and not raining across town; it might be raining now but stop raining even as you finish this sentence; but it is impossible for it to be both raining and not raining in the same place and at the same time. So, whatever the world is like, it is not both raining here and not raining here. It is a \define{necessary falsehood}.

\newglossaryentry{necessary falsehood}
{
name={necessary falsehood},
description={A sentence that is false in every case}
}

Something might \emph{always} be true and still be contingent. For instance, if there never were a time when the universe contained fewer than seven things, then the sentence `At least seven things exist' would always be true. Yet the sentence is contingent: the world could have been much, much smaller than it is, and then the sentence would have been false.

\section{Necessary equivalence}

We can also ask about the logical relations \emph{between} two sentences. For example:
\begin{compactlist}
\item[] John went to the store after he washed the dishes.
\item[] John washed the dishes before he went to the store.
\end{compactlist}
These two sentences are both contingent, since John might not have gone to the store or washed dishes at all. Yet they must have the same truth-value. If either of the sentences is true, then they both are; if either of the sentences is false, then they both are. When two sentences have the same truth value in every case, we say that they are \define{necessarily equivalent}.

\newglossaryentry{necessary equivalence}
{
name={necessary equivalence},
text={necessarily equivalent},
description={A property held by a pair of sentences that, in every case, are either both true or both false}
}


\section*{Summary of logical notions}

\factoidbox{
\begin{factoiditems}
\item An argument is \define{valid} if there is no case where the premises are all true and the conclusion is not; it is \define{invalid} otherwise.

\item A \define{necessary truth} is a sentence that is true in every case.

\item A \define{necessary falsehood} is a sentence that is false in every case.

\item A \define{contingent sentence} is neither a necessary truth nor a necessary falsehood; a sentence that is true in some case and false in some other case.

\item Two sentences are \define{necessarily equivalent} if, in every case, they are both true or both false.

\item A collection of sentences is \define{jointly possible} if there is a case where they are all true together; it is \define{jointly impossible} otherwise.
\end{factoiditems}
}

\practiceproblems
\problempart
\label{pr.EnglishTautology2}
For each of the following: Is it a necessary truth, a necessary falsehood, or contingent?
\begin{compactlist}
\item Caesar crossed the Rubicon.
\item Someone once crossed the Rubicon.
\item No one has ever crossed the Rubicon.
\item If Caesar crossed the Rubicon, then someone has.
\item Even though Caesar crossed the Rubicon, no one has ever crossed the Rubicon.
\item If anyone has ever crossed the Rubicon, it was Caesar.
\end{compactlist}

\problempart
For each of the following: Is it a necessary truth, a necessary falsehood, or contingent?
\begin{compactlist}
\item Elephants dissolve in water.
\item Wood is a light, durable substance useful for building things.
\item If wood were a good building material, it would be useful for building things.
\item I live in a three-story building that is two stories tall.
\item If gerbils were mammals, they would nurse their young.
\end{compactlist}

\problempart Which of the following pairs of sentences are necessarily  equivalent? 

\begin{compactlist}
\item Elephants dissolve in water.	\\
	If you put an elephant in water, it will disintegrate.
\item All mammals dissolve in water.\\		
	If you put an elephant in water, it will disintegrate.
\item George Bush was the 43rd president. \\
	 Barack Obama is the 44th president.
\item Barack Obama is the 44th president. \\
	  Barack Obama was president immediately after the 43rd president.
\item Elephants dissolve in water. 	\\	
	All mammals dissolve in water.
\end{compactlist}
\problempart Which of the following pairs of sentences are necessarily equivalent? 

\begin{compactlist}
\item  Thelonious Monk played piano.	\\
	John Coltrane played tenor sax.
\item  Thelonious Monk played gigs with John Coltrane.	\\
	John Coltrane played gigs with Thelonious Monk.
\item  All professional piano players have big hands.	\\
	Piano player Bud Powell had big hands.
\item  Bud Powell suffered from severe mental illness.	 \\
	All piano players suffer from severe mental illness.
\item John Coltrane was deeply religious.	 \\
John Coltrane viewed music as an expression of spirituality.
\end{compactlist}

\noindent \problempart Consider the following sentences: 
\begin{compactlist}%[label=(\alph*)]
\item[G1.] \label{itm:at_least_four}There are at least four giraffes at the wild animal park.
\item[G2.] \label{itm:exactly_seven} There are exactly seven gorillas at the wild animal park.
\item[G3.] \label{itm:not_more_than_two} There are not more than two Martians at the wild animal park.
\item[G4.] \label{itm:martians} Every giraffe at the wild animal park is a Martian.
\end{compactlist}

Now consider each of the following collections of sentences. Which are jointly possible? Which are jointly impossible?
\begin{compactlist}
\item Sentences G2, G3, and G4
\item Sentences G1, G3, and G4
\item Sentences G1, G2, and G4
\item Sentences G1, G2, and G3
\end{compactlist}

\problempart Consider the following sentences.
\begin{compactlist}%[label=(\alph*)]
\item[M1.] \label{itm:allmortal} All people are mortal.
\item[M2.] \label{itm:socperson} Socrates is a person.
\item[M3.] \label{itm:socnotdie} Socrates will never die.
\item[M4.] \label{itm:socmortal} Socrates is mortal.
\end{compactlist}
Which combinations of sentences are jointly possible? Mark each ``possible'' or ``impossible.''
\begin{compactlist}
\item Sentences M1, M2, and M3
\item Sentences M2, M3, and M4
\item Sentences M2 and M3
\item Sentences M1 and M4
\item Sentences M1, M2, M3, and M4
\end{compactlist}

\problempart
\label{pr.EnglishCombinations2}
Which of the following are possible? For each, if it is possible, give an example. If it is not possible, explain why.
\begin{compactlist}
\item A valid argument that has one false premise and one true premise

\item A valid argument that has a false conclusion

\item A valid argument, the conclusion of which is a necessary falsehood

\item An invalid argument, the conclusion of which is a necessary truth

\item A necessary truth that is contingent

\item Two necessarily equivalent sentences, both of which are necessary truths

\item Two necessarily equivalent sentences, one of which is a necessary truth and one of which is contingent

\item Two necessarily equivalent sentences that together are jointly impossible

\item A jointly possible collection of sentences that contains a necessary falsehood

\item A jointly impossible set of sentences that contains a necessary truth
\end{compactlist}

\problempart
Which of the following are possible? For each, if it is possible, give an example. If it is not possible, explain why.

\begin{compactlist}
\item A valid argument, whose premises are all necessary truths, and whose conclusion is contingent
\item A valid argument with true premises and a false conclusion
\item A jointly possible collection of sentences that contains two sentences that are not necessarily equivalent
\item A jointly possible collection of sentences, all of which are contingent
\item A false necessary truth
\item A valid argument with false premises
\item A necessarily equivalent pair of sentences that are not jointly possible
\item A necessary truth that is also a necessary falsehood
\item A jointly possible collection of sentences that are all necessary falsehoods
\end{compactlist}
